% \documentclass{fp-slides}
\documentclass[10pt,portrait]{beamer}
\usetheme{Madrid}
% \usetheme{Warsaw}
\institute{O.O.Bogomolets National Medical University}
% \ifx\pdfoutput\undefined
% % we are running LaTeX, not pdflatex
% \usepackage{graphicx}
% \else
% % we are running pdflatex, so convert .eps files to .pdf
% \usepackage[pdftex]{graphicx}
% \usepackage{epstopdf}
% \fi 
\usepackage{graphics}

\linespread{1.1}
% \date{}
\title{Introdution to medical informatics. Data and information.}
\author[Shiray Andrey]{Shiray Andrey}

\begin{document}
 \frame{\titlepage}

 \frame{
   \frametitle{Outline}
 
   \begin{itemize}
   \item \Large Course overview

   \item Medical informatics. Why? What for?

\item Fundamental role of informatics in modern science.
   
  \item Information. Data. Knowledge.


   \end{itemize}
 }

\frame{
\frametitle{Course overview}
\begin{itemize}
\item \Large Information. Data. Knowledge
\item Biostatistics
\item Formal Logic. Expert and decision support systems in medicine
\item Evidence-based medicine
  \end{itemize}
}

\frame{
\frametitle{Medical informatics}
\Huge
\textbf{Medical informatics}  -- the application of informatics to the field of biology and medicine/healthcare. 
% (also called health informatics, health care informatics, healthcare informatics, clinical informatics, or biomedical informatics)
}

\frame{
\frametitle{Informatics, Computers science}
\begin{itemize}
\item \Large \textbf{Informatics} -- the study of the processing, management, and retrieval of information
\item \textbf{Computers science} -- is the study of the theoretical foundations of information and computation and of practical techniques for their implementation and application in computer systems
\item \textbf{Cybernetics} is the interdisciplinary study of the structure of regulatory systems.
\end{itemize}
}

\frame{
\frametitle{Medical informatics. Why? What for?}
\Large Medical informatics and borrows concepts/ideas from:
\begin{itemize}
 \item Medicine
\item Informatics
\item Logic
\item Statistics
\end{itemize}
}

\frame{
\frametitle{Fundamental role of informatics in modern science}
\begin{itemize}
 \item \Large Any object in physical reality can be described with some information

\end{itemize}

}

\frame{
\frametitle{Fundamental role of informatics in modern science}
\begin{itemize}
 \item \Large Any object in physical reality can be described with some information
\item Any physical process can be described as informational process
\item Physical laws $\sim  $ program/algorithm
\end{itemize}

}

\frame{
\frametitle{Fundamental role of informatics in modern science}
\begin{itemize}
 \item \Large Any object in physical reality can be described with some information
\item Any physical process can be described as informational process
\item Physical laws $\sim  $ program/algorithm
\item Universe is a gigantic computer!
\begin{itemize}
\item \footnotesize Seth Lloyd, Programming the Universe
\item \footnotesize Roger Penrose, The Emperor's New Mind
\end{itemize}
\end{itemize}

}



\frame{
\frametitle{Fundamental role of informatics in modern society}

That's why informatics is an essential part of science disciplines in any University.


You will need basic informatics skills in modern society:

\begin{itemize}
\item We live in so called information society. \footnotesize(or at least tend to develop it)
 \item You may need to analyze systems/products
\item You may need apply informatics skills  in your research projects
\item You may need to implement prototypes
\item You may need to oversee programmers in your future job
\end{itemize}

}

\frame{
\frametitle{Fundamental concepts of informatics}
\begin{itemize}
 \item \Large Information 
\item Computation
\item Data
\item Knowledge
\end{itemize}

}

\frame{
\frametitle{Information}
\Huge \textbf{Information} -- primitive notion\footnote{fundamental undefined concept; primitive, undefined, term} of modern science.

\vspace{5mm}
\normalsize So  it is not defined formally in terms of previously defined concepts, but is only motivated informally, usually by an appeal to intuition and everyday experience
}

\frame{
\frametitle{Properties}
\Large Properties of information:

\begin{itemize}
 \item \textbf{Objectivity} -- information doesn't depend upon storage method or your opinion
\item \textbf{Authenticity} -- how trustworsy it is
\item \textbf{Accuracy} 
\item \textbf{Actual}
 
\end{itemize}

}


\frame{
\frametitle{Measures of Information. Shannon's  measure}
In \textit{information theory}, \textbf{entropy} is a measure of the uncertainty associated with a random variable(in our case received message)


\vspace{5mm}
Shannon's  measure\footnote{also caled self-information }:
$I(x) = \log (\frac{1}{P(x)}) = - \log(P(x)) $

\vspace{5mm}
Example:


On tossing a coin, the chance of 'tail' is $0.5$. When it is proclaimed that indeed 'tail' occurred, this amounts to
$I('tail') = log_2 (\frac{1}{0.5}) = \log_2 2 = 1$ bits of information.
}

%%%%%%%%%% Break


\frame{
\frametitle{Data}
\textbf{Data} refers to qualitative or quantitative attributes of a variable or set of variables.

\begin{itemize}
\item For example -- results of measurements: $(35.7, 36.7, 42, 38, 39.4, 37.5, 40.1)$
\item Raw data is the lowest level of abstraction. Raw data on its own carries no meaning.
\end{itemize}


}

\frame{
\frametitle{Data}
\textbf{Data} refers to qualitative or quantitative attributes of a variable or set of variables.

\begin{itemize}
\item For example -- results of measurements: $(35.7, 36.7, 42, 38, 39.4, 37.5, 40.1)$
\item Raw data is the lowest level of abstraction. Raw data on its own carries no meaning.
\item Raw data needs some way of interpretation.
\end{itemize}


}

\frame{
\frametitle{Data}
\textbf{Data} refers to qualitative or quantitative attributes of a variable or set of variables, stored in some way.

\begin{itemize}
\item For example -- results of measurements: $(35.7, 36.7, 42, 38, 39.4, 37.5, 40.1)$
\item Raw data is the lowest level of abstraction. Raw data on its own carries no meaning.
\item Raw data must be interpreted and take on a meaning
\item Example: $(35.7, 36.7, 42, 38, 39.4, 37.5, 40.1)$ -- results of patients thermometrics
\end{itemize}


}

\frame{
\frametitle{Knowledge}
\textbf{Knowledge} is a familiarity with someone or something, that can include facts/data, descriptions, information, and/or skills acquired through experience or education.


\vspace{5mm}


Main difference between knowledge and data: knowledge has dynamic nature by itself, it can change via learning. This the highest level of information abstraction.
}
\frame{
\frametitle{Computation}
\begin{itemize}

\item \Large \textbf{Computation} -- any type of information processing.

\end{itemize}
}
\frame{
\frametitle{Computation}
\begin{itemize}

\item \Large \textbf{Computation} -- any type of information processing.


\item A computation can be seen as a purely physical phenomenon occurring inside a closed physical system called a computer.

\end{itemize}
}
\frame{
\frametitle{Computation}
\begin{itemize}

\item \Large \textbf{Computation} -- any type of information processing.


\item A computation can be seen as a purely physical phenomenon occurring inside a closed physical system called a computer.


\item To describe a computation we need some formal way to state rules for information processing - \textbf{algorithm}.
\end{itemize}
}

\frame{
\frametitle{Algorithm}
\Large \textbf{Algorithm}  is an effective method expressed as a finite list of well-defined instructions for computation.

\vspace{5mm}
Starting from an initial state and initial input (perhaps empty), the instructions describe a computation that, when executed, will proceed through a finite number of well-defined successive states,using finite amount of resources, eventually producing "output"
}

\frame{
\frametitle{Properties of Algorithms}
\begin{itemize}
\item  \textbf{Finiteness} -- An algorithm must always terminate after a finite number of steps
\item  \textbf{Definiteness} -- Each step of an algorithm must be precisely defined; the actions to be carried out must be rigorously and unambiguously specified for each case
% \item Discrete
\item \textbf{Discrete} --  moving in separate(discrete) steps: 1, 2, 3, ...
\item \textbf{Input}
\item \textbf{Output}
% \item 
\end{itemize}

}
% \frame{
% \frametitle{Kolmogorov complexity. Human Genome.}
% In algorithmic information theory, the \textbf{Kolmogorov complexity} of an object, such as a piece of text, is a measure of the computational resources needed to specify the object.
% 
% }
% \frame{
% \frametitle{Flowcharts}
%  \begin{center}
%  \begin{tabular}{ | m{2cm} | m{2cm} | m{5cm} |}
% % \noalign{\smallskip}
% % Name & Drawing & Function\\ \hline
% % Start and end symbols & \includegraphics[scale=0.5]{ss} & signaling the start or end of a process\\  \hline
% % Generic processing steps & \includegraphics[scale=0.5]{proc} & Examples: "Add 1 to X"; "replace identified part"; "save changes" or simila \\   \hline
% % Conditional or decision & \includegraphics[scale=0.5]{if} &  showing where a decision is necessary, commonly a Yes/No question or True/False test 
% % 
%  \end{tabular}
%  \end{center}
% 
% }

% \frame{
% \frametitle{Entropy. Complex systems.}
%  \textbf{Entropy} -- amount of disorder in a system
% 
% 
% }
% 
% \frame{
% \frametitle{Information Ontology}
% }





\end{document}
